% Source: sources/returns/S6_04_MAPS_MANIFOLDS/sections/manifolds.tex:965–1148
\paragraph{Theorem D.2.5 (Proved: arbitrary multiplicities in the original configuration).}
Consider the $IV^*,III,I_1$ configuration with its rank-two Mumford
filling, and the logarithmic data of
\cite[\S2.1, (2.11)--(2.12)]{manViaclovskyLogs26}:
\[
 m_0=3k,\qquad m_1=4\ell,\qquad
 v_i=r_i\delta+n_i u_i,\qquad
 \gcd(n_i,m_i)=1,\qquad k,\ell>0.
\]
Here $r_i,n_i$ are arbitrary integers, not residue classes. Put
$P=m_1n_0+m_0n_1$. There is no upper bound on these parameters.
If $P\ne0$, the resulting smooth compact threefold has
\[
 \pi_1(X)=\mathbb Z/|P|,\qquad
 H_j(X;\mathbb Z)=
 \begin{cases}
 \mathbb Z&j=0,6,\\
 \mathbb Z/|P|&j=1,2,3,4,\\
 0&j=5.
 \end{cases}
\]
If $P=0$, then $\pi_1(X)=\mathbb Z$ and all integral homology
is free, with Betti numbers $(1,1,2,2,2,1,1)$.
In particular, in this admissible family $X$ is simply connected
exactly when $|P|=1$, and then it is diffeomorphic to $S^6$.
Such sphere data exist exactly when $\gcd(3k,4\ell)=1$.
This proves a parametric table for every multiplicity and every
integer twist in this family; it does not impose the determinant-one
hypothesis when computing the other manifolds.
\begin{proof}
We give the integral reduction explicitly, including the step that
extends beyond the simply connected hypothesis in
\cite[Propositions 4.1, 4.3, 4.5]{manViaclovskyLogs26}.
In the integral basis $(\delta,a,b,d)$, the three homological
monodromies are
\[
 A_0=\begin{pmatrix}1&-1&0&2\\0&-1&-1&6\\0&1&0&0\\0&0&0&1\end{pmatrix},\quad
 A_1=\begin{pmatrix}1&0&1&-3\\0&0&1&0\\0&-1&0&6\\0&0&0&1\end{pmatrix},\quad
 A_\infty=\begin{pmatrix}1&0&0&1\\0&1&1&0\\0&0&1&0\\0&0&0&1\end{pmatrix}.
\]
They satisfy $A_0A_1A_\infty=I$, $A_0^3=A_1^4=I$,
$u_0=d+2a+2b$, $u_1=d+3a+3b$ and
$\Lambda^{A_i}=\mathbb Z\delta\oplus\mathbb Zu_i$ for $i=0,1$.
The good reduction followed by the degree-$k$ or degree-$\ell$
base change remains smooth. The affine generator has translation
$v_i/m_i$ and rotates the base by $e^{2\pi i/m_i}$.
The invariant covector $d^\vee$ changes by $n_i/m_i$;
$\gcd(n_i,m_i)=1$ therefore proves that every nonidentity power
is free even on the central torus. Quotienting gives a smooth
filling retracting onto a bielliptic surface $F_i$.
The logarithm translation used in the punctured-disc gluing has
normal-circle relation $x_i^{m_i}=v_i$. These are precisely the
charts and gluing of Proposition 2.3 of the cited paper; the
freeness argument uses no relation between the two multiplicities.

The Mumford filling kills $\delta,a,x_\infty$ in the fundamental
group. The conjugation relation $x_1ax_1^{-1}=b^{-1}$ then kills
$b$, and $x_1=x_0^{-1}$. Thus the full group is already abelian,
generated by $d$ and $x=x_0$, with relation matrix
\[
 R=\begin{pmatrix}-n_0&m_0\\-n_1&-m_1\end{pmatrix}.
\]
Its entries have greatest common divisor one because
$\gcd(n_0,m_0)=1$, and $\det R=P$. Its Smith form is therefore
$(1,|P|)$ when $P\ne0$, or $(1,0)$ when $P=0$.
This is a computation of the group, not just its abelianization.

Let $V$ be the inverse image of the thrice-punctured sphere, and
$M_i=\partial N_i$ the boundary of each filling. All $H_1(M_i)$
are free of rank three: the individual coinvariants give
$b=-a,\delta=-3a$ at $0$, and $b=-a,\delta=-2a$ at $1$;
the cusp has primitive rank-two vanishing lattice.
The common coinvariants give $H_1(V)=\mathbb Z^3$.
Independently of $P$, the finite central surfaces have
\[
 H_1(F_0)=\mathbb Z\langle a,d,g_0\rangle/
   \langle3r_0a-n_0d+m_0g_0\rangle,
\quad
 H_1(F_1)=\mathbb Z\langle a,d,g_1\rangle/
   \langle2r_1a-n_1d+m_1g_1\rangle.
\]
These are free of rank two. Choose
$m_i\alpha_i+n_i\beta_i=1$ and put
$z_i=\alpha_i d+\beta_i g_i$; $(a,z_i)$ is an integral basis.
Also $H_1(F_\infty)=\mathbb Z^2$.
Consequently the degree-one Mayer--Vietoris map is a map between
free rank-nine groups with cokernel $H_1(X)$.
If $P\ne0$ its kernel is zero, even though its cokernel may be
nonzero. If $P=0$ its kernel is free of rank one.
This replaces the simply connected assumption in the cited
Proposition 4.1 by the actual rank computation.

The map $H_2(M_i)\to H_2(N_i)$ is onto for every filling.
Indeed its cokernel maps to
$H_2(N_i,M_i)\simeq H^4(F_i)=\mathbb Z$ by intersection
with $F_i$; the multiple $m_iF_i$ is principal, so that
intersection is zero. By the pair sequence, its kernel consists
of the normal-circle tori of classes in $H^3(F_i)$.
The homology Serre sequence over the graph underlying the base gives
\[
 0\longrightarrow (\bigwedge^2\Lambda)_{A_0,A_1}
 \longrightarrow H_2(V)\longrightarrow
 \ker\bigl[(A_0-I)\ (A_1-I)\bigr]\longrightarrow0.
\]
The left group is $\mathbb Z[a\wedge b]$: the exterior relations
give $\delta\wedge a=\delta\wedge b=a\wedge d=b\wedge d=0$
and $\delta\wedge d=-6a\wedge b$.
The invariant form
$\xi=a^\vee\wedge b^\vee+6d^\vee\wedge\delta^\vee$
evaluates to one on $a\wedge b$, proving primitivity.
The right group is free of rank five.
The two cusp normal tori are the sweeps of $\delta,a$.
In the graph kernel they are
$(-\delta,-\delta)$ and $(b,-a)$; quotienting by them leaves
the integral basis $(\delta,0),(u_0,0),(0,u_1)$.
For verification, the equation in the $A_0$ coinvariants forces
$b^\vee(v_1)=3d^\vee(v_1)$; subtracting the appropriate cusp
vector then puts $v_1$ in $\mathbb Zu_1$ and $v_0$ in
$\mathbb Z\delta\oplus\mathbb Zu_0$. Uniqueness follows
by comparing the second components.
Thus, after the cusp relations, the target for the remaining
normal tori is a free rank-four group with basis $(s,t,t_0,t_1)$,
where $s=[a\wedge b]$ and the last three classes are these sweeps.

For a loop $\lambda$ in the finite covering torus, the normal
torus cuts the reference fibre in
$\mathcal N_i\lambda$, where
$\mathcal N_i=\sum_{j=0}^{m_i-1}A_i^j$, and its $s$ coordinate
is $-\xi(\lambda,v_i)$. The normal torus of $g_i$ has
fibre cut $v_i$ and $s$ coordinate zero; both follow by applying
the homomorphisms to $m_i g_i=v_i$ and cancelling in free groups.
The norms are
\[
 \mathcal N_0a=-k\delta,\quad \mathcal N_0d=m_0u_0,
 \qquad
 \mathcal N_1a=-2\ell\delta,\quad \mathcal N_1d=m_1u_1.
\]
Moreover $\xi(a,u_0)=2$, $\xi(a,u_1)=3$ and
$\xi(d,\delta)=6$. Since $(0,\delta)=-(\delta,0)$ modulo
the cusp relations, the normal-torus map in the bases
$(a,z_0,a,z_1)$ and $(s,t,t_0,t_1)$ is
\[
 B=\begin{pmatrix}
 -2n_0&-6r_0\alpha_0&-3n_1&-6r_1\alpha_1\\
 -k&r_0\beta_0&2\ell&-r_1\beta_1\\
 0&1&0&0\\
 0&0&0&1
 \end{pmatrix}.
\]
Removing the two unit pivots reduces its cokernel to that of
\[
 C=\begin{pmatrix}-2n_0&-3n_1\\-k&2\ell\end{pmatrix},
 \qquad \det C=-P.
\]
The entries of $C$ have greatest common divisor one: any common
divisor divides $\gcd(k,2n_0)$ and hence divides $2$, whereas
$n_1$ is odd because $\gcd(n_1,4\ell)=1$.
Thus $\operatorname{coker}B=\mathbb Z/|P|$ for $P\ne0$
and $\mathbb Z$ for $P=0$. Mayer--Vietoris gives $H_2(X)$
equal to this cokernel for $P\ne0$; for $P=0$ it is an extension
of $\mathbb Z$ by $\mathbb Z$, hence $\mathbb Z^2$.

Every fibre except the rank-two cusp has Euler number zero, so
$\chi(X)=2$. Integral Poincare duality and the universal
coefficient theorem give $H_5\simeq\operatorname{Hom}(H_1,\mathbb Z)$,
$H_4\simeq\operatorname{Hom}(H_2,\mathbb Z)
\oplus\operatorname{Ext}(H_1,\mathbb Z)$, and
$\operatorname{Tor}H_3\simeq\operatorname{Tor}H_2$.
The Euler characteristic determines the remaining rank, giving
all stated groups. For $|P|=1$, sphere recognition follows as in
D.3.2 from Hurewicz, Whitehead and $\Theta_6=0$.
Finally the Bezout equation $m_1n_0+m_0n_1=\pm1$ is solvable
exactly when $\gcd(m_0,m_1)=1$, and it automatically implies both
local admissibility conditions.

The unrestricted exact implementation is the script
\texttt{manifolds\_unbounded\_logs.py} in \texttt{code/}. It symbolically checks
the determinant and norm formulas and actually computes the Smith
forms on the stated finite validation grid; its all-parameter
correctness is the integral argument above. For example,
$(m_0,m_1;n_0,n_1)=(9,4;-2,1)$ gives $P=1$ and $S^6$;
$(9,4;1,1)$ gives all four torsion groups $\mathbb Z/13$;
$(12,12;1,-1)$ gives the stated infinite-cyclic case.
\end{proof}
